% !TEX root = ../frode-bibelen.tex
% ----------------------------------------------------------------------------
%  Frode-bibelen -- 11 Struct, pekere og memory
%  Inkluderes fra frode-bibelen.tex (kompiler alltid hovedfilen).
% ----------------------------------------------------------------------------

% ================================================== 11 STRUCT/PEKERE ========
\kapittel{Struct, pekere og memory}

\subsection{Pekere på én side}
\begin{center}
\begin{tabular}{@{}L{3.6cm}L{12.0cm}@{}}
\toprule
\cc{int* p;} & Lager en peker som kan peke på en \cc{int}. \\
\cc{p = \&x;} & \cc{\&x} = adressen til \cc{x}. Nå peker \cc{p} på \cc{x}. \\
\cc{*p = 7;} & \cc{*p} = verdien der \cc{p} peker (her: setter \cc{x} til 7). \\
\cc{t == \&t[0]} & Navnet på en array \textbf{er} adressen til første element. \\
\cc{*(t+i) == t[i]} & Pekeraritmetikk: \cc{t+i} peker \cc{i} plasser lenger ut. \\
\cc{p->navn} & Medlem via peker til struct. Samme som \cc{(*p).navn}. \\
\cc{s.navn} & Medlem i en struct-\emph{variabel} (ikke peker). \\
\cc{NULL} & «Peker ikke på noe». \cc{if (p)} = \cc{if (p != NULL)}. \\
\bottomrule
\end{tabular}
\end{center}

\subsection{\texorpdfstring{\cc{malloc}}{malloc} og \texorpdfstring{\cc{free}}{free}}
\begin{kode}
struct Kunde* ny = (struct Kunde*) malloc(sizeof(struct Kunde));   //  Én struct.
char* t  = (char*) malloc((strlen(buffer) + 1) * sizeof(char));    //  Én tekst.
int*  a  = (int*)  malloc(n * sizeof(int));                        //  Array med n int.

free(kunde->navn);       //  1) Først det som er malloc-et INNI structen,
free(kunde);             //  2) så selve structen.   (Omvendt rekkefølge = minnelekkasje.)
\end{kode}

\subsection{Array \emph{av} structer vs.\ array \emph{med pekere til} structer}
Begge varianter har vært på eksamen (S26 hadde begge i samme program).
\begin{center}
\begin{tabular}{@{}L{3.0cm}L{6.1cm}L{6.5cm}@{}}
\toprule
 & \textbf{Array av structer} & \textbf{Array med pekere (vanligst)} \\
\midrule
Deklarasjon & \cc{struct Glass gGlass[MAXGLASS];} & \cc{struct Kunde* gKunder[MAXKUNDER];} \\
Tilgang & \cc{gGlass[i].bredde} & \cc{gKunder[i]->navn} \\
Sende til funksjon & \cc{glassSkriv(gGlass[i])} (kopi) eller \newline \cc{glassLes(\&gGlass[i])} (original) & \cc{kundeSkrivData(gKunder[i])} \newline (er allerede en peker) \\
Opprette ny & ingenting -- plassen finnes & \cc{gKunder[n] = (struct Kunde*)} \newline \cc{\ \ malloc(sizeof(struct Kunde));} \\
Slette & ikke aktuelt / flytt innhold & \cc{free} inni, \cc{free} structen, flytt siste peker inn \\
\bottomrule
\end{tabular}
\end{center}

\subsection{Typiske ting \emph{inni} en struct}
\begin{kode}
struct Fugl {
  char* navn;                   //  Dynamisk tekst.
  int   antObs;                 //  Teller  +  array som brukes 0..antObs-1
  int   observasjoner[365];     //    (som gKunder/gAntallKunder, men inni structen).
  bool  dag[31];                //  Ja/nei for hver dag  ->  dag[d-1] for dag nr. d
  int   bankNr[2];              //  INDEKSER inn i en annen global array (relasjon):
};                              //    gBanker[ kunde->bankNr[0] ]->navn

struct Kunde {
  int   glass[MAXGLASS];        //  Antall bestilt av glass nr. i (S26): én teller
};                              //    per element i en annen global array.
\end{kode}

\begin{tips}[{Kopi av struct er «grunn»}]
\lstinline|*s1 = *s2;| kopierer alle medlemmene -- men for en \cc{char*} kopieres bare \emph{pekeren}.
Begge structene peker da på samme tekst. (\cc{char navn[30]} inni structen blir derimot kopiert.)
\end{tips}

\subsection{Datoer og klokkeslett som heltall}\label{sec:dato}
Frode lagrer ofte dato som \cc{MMDD} (H22, H23, S26) og klokkeslett som \cc{TTMM} (H22):
\begin{kode}
int dato = 1224;                        //  24. desember
int mnd  = dato / 100;                  //  12
int dag  = dato % 100;                  //  24
printf("%02i/%02i", dag, mnd);          //  24/12

int dato2 = mnd * 100 + dag;            //  Tilbake til MMDD.
if (dato1 < dato2)  ...                 //  MMDD kan sammenlignes direkte (samme år).

int klokken = 1430;                     //  14:30
int minutterEtterMidnatt = (klokken / 100) * 60 + klokken % 100;   //  870

const int DAGERIMND[12] = { 31,28,31,30,31,30,31,31,30,31,30,31 };
int dagNr = dag;                                    //  Dag nummer i året (1..365):
for (int m = 0;  m < mnd-1;  m++)  dagNr += DAGERIMND[m];
\end{kode}
